\section{What bulk observables hide, and what fails beyond the additive kernel}
\label{sec:boundaries}

Count neutrality has an exact algebraic classification at fixed mass. It
explains why constant count, mass, and mean size can coexist with rapid
broadening. A second observable removes this universal ambiguity. Beyond
the additive kernel, the same classification also leads to a finite-count
nonstationarity theorem, while fractional-moment monotonicity can fail.

\subsection{Universal count neutrality and observable rigidity}

In this subsection $K:E^2\to[0,\infty)$ is measurable, symmetric and finite
pointwise, and $S:E\to[0,\infty)$ is measurable and finite pointwise. The
expected daughter measure $b_x$ satisfies \eqref{eq:daughters}. For a
finitely supported preparation $v$, define its count vector field by
\[
 \mathcal R(v)=\int S(x)v(\dd x)-\frac12\iint K(x,y)v(\dd x)v(\dd y).
\]
This definition concerns initial rates and needs no existence theorem for
the general kernel. Fix a prescribed material mass $m>0$. The preparation
class below fixes this mass while allowing count to vary.

\begin{theorem}[The complete count-neutral class]\label{thm:neutral-class}
The identity $\mathcal R(v)=0$ holds for every nonnegative finitely
supported $v$ with $M_1(v)=m$ if and only if there is a nonnegative
measurable function $d$, finite pointwise, such that
\begin{equation}
 K(x,y)=xd(y)+yd(x),\qquad S(x)=md(x).
 \label{eq:neutral-class}
\end{equation}
For this class, any mass-conserving trajectory at mass $m$ whose weak count
balance has locally integrable event rates has constant count.
\end{theorem}

\begin{proof}
The preparation $(m/x)\delta_x$ gives
$S(x)=mK(x,x)/(2x)$. Define $d(x)=S(x)/m$. For $x\ne y$, substitute
$v=(m/(2x))\delta_x+(m/(2y))\delta_y$ and use the diagonal identities.
After cancellation the result is $K(x,y)=xd(y)+yd(x)$; the same relation
already holds on the diagonal. Conversely, for \eqref{eq:neutral-class},
\begin{equation}
 \mathcal R(v)=(m-M_1(v))\int d(x)v(\dd x).
 \label{eq:neutral-count-balance}
\end{equation}
This proves the equivalence. Along a trajectory with the stated event
integrability, the weak count balance integrates the right-hand side of
\eqref{eq:neutral-count-balance}, which vanishes by mass conservation.
\end{proof}

The full-time assertion uses the weak equation and conserved mass, rather
than inferring a trajectory from an initial derivative. The fixed-mass
preparation class is essential: a class fixing both mass and count admits
only one monodisperse size and cannot support the proof below.

\begin{proposition}[An additional moment forces zero kinetics]
\label{prop:observable-rigidity}
Fix any $0<p<1$. If the count vector field vanishes for every finitely
supported preparation of mass $m$, and the $p$th-moment vector field
vanishes for every monodisperse preparation of mass $m$, then
$K=S=0$ pointwise. The same conclusion holds if the second moment replaces
the $p$th moment. In particular one may use the half moment.
\end{proposition}

\begin{proof}
By \cref{thm:neutral-class}, the rates have form \eqref{eq:neutral-class}
with $d\geq0$. At $v=(m/x)\delta_x$ the coagulation contribution to the
$p$th-moment vector field is
$m^2d(x)x^{p-1}(2^p-2)$. Jensen's inequality bounds fragmentation by
$m^2d(x)x^{p-1}(2^{1-p}-1)$. Therefore their sum is at most
\begin{equation}
 -\kappa_p m^2d(x)x^{p-1}.
 \label{eq:rigidity-fractional}
\end{equation}
If this vector field is zero, then $d(x)=0$ for every $x$.
For the second moment, Jensen gives
$\int z^2b_x(\dd z)\geq x^2/2$, so the corresponding vector field is
at least $3m^2xd(x)/2$. Its universal vanishing again forces $d=0$.
Every displayed integral is finite for these monodisperse preparations
because daughters lie below their parent.
\end{proof}

This is a result about a class of preparations at fixed mass, not
identification of arbitrary rates from a single observed trajectory.
Only monodisperse preparations enter the moment hypothesis, and the proof
of \cref{thm:neutral-class} uses only preparations with at most two atoms,
so the count hypothesis can be weakened to that class as well. Nonnegative
physical rates are used in the conclusion from
\eqref{eq:rigidity-fractional}.

\subsection{Broadening and entropy production with fixed bulk outputs}

Return to the additive member $d\equiv\lambda>0$, with critical selection
$b=\lambda m$. The daughter law may now depend measurably on time and size,
as in \cref{ass:rates}. In the finite-second-moment solution class already
constructed, second-moment and entropy balances require no extra assumptions.
Use $\mathcal H(t)=\int x\log x\,n_t(\dd x)$ for the entropy moment, to
avoid confusing it with the half moment $H$ used earlier.

\begin{proposition}[Two unbounded observables]\label{prop:broadening-entropy}
For the critical solution of \cref{thm:existence}, $M_2$ and $\mathcal H$
are finite and locally absolutely continuous, and almost everywhere
\begin{align}
 \frac32\lambda mM_2(t)&\leq M_2'(t)\leq2\lambda mM_2(t),
 \label{eq:broadening}\\
 \mathcal H'(t)&\geq\lambda m^2\log2.
 \label{eq:entropy-production}
\end{align}
In particular,
\begin{equation}
 M_2(0)e^{3\lambda mt/2}\leq M_2(t)\leq M_2(0)e^{2\lambda mt},
 \qquad
 \mathcal H(t)\geq\mathcal H(0)+\lambda m^2t\log2.
 \label{eq:broadening-integrated}
\end{equation}
Equal splitting gives exactly $M_2(t)=M_2(0)e^{3\lambda mt/2}$.
The entropy production constant is sharp as a uniform instantaneous bound:
monodisperse data with equal splitting attain it at time zero.
\end{proposition}

\begin{proof}
The unbounded-test justification is given in \cref{app:unbounded}; it
establishes the complete integrated weak balances using only locally
bounded $M_2$. The coagulation contribution to $M_2'$ is
$2\lambda mM_2$. By Jensen and $z<x$,
\[
 x^2/2\leq\int z^2b_{t,x}(\dd z)\leq x^2,
\]
so the fragmentation contribution lies between $-bM_2/2$ and zero.
This proves \eqref{eq:broadening} and its equal-split identity.

For $f(x)=x\log x$, put $u=x/(x+y)$ and
$h(u)=-u\log u-(1-u)\log(1-u)$. Then
$\Delta f(x,y)=(x+y)h(u)$. Divide \eqref{eq:pair} by $1-p$ and let
$p\uparrow1$ at fixed positive $x,y$. Differentiability of the scalar powers
gives
\begin{equation}
 (x+y)\Delta f(x,y)\geq4(\log2)xy,
 \quad\text{equivalently}\quad h(u)\geq4(\log2)u(1-u).
 \label{eq:entropy-pair}
\end{equation}
Estimates of this subadditivity type for concave tests are standard in
the coagulation--fragmentation literature; see
\citet{EscobedoMischlerRodriguezRicard2005}.
Thus coagulation produces at least $2\lambda m^2\log2$. Jensen's inequality
for $b_{t,x}/2$ and the increase of $f(z)/z=\log z$ give
\[
 -x\log2\leq\int f(z)b_{t,x}(\dd z)-f(x)\leq0.
\]
Fragmentation loses at most $bm\log2$, proving
\eqref{eq:entropy-production}. Integrating proves
\eqref{eq:broadening-integrated}. Both scalar inequalities are equalities
at a monodisperse equal-split initial state. The right derivative there
is justified at the end of \cref{app:unbounded}, so sharpness is a statement
about actual solutions as well as their vector fields.
\end{proof}

Expected daughter count two and expected daughter mass $x$ suffice for
these bounds; an eventwise binary law is not needed. For any fixed initial
measure with $0<M_2(0)<\infty$, compare the equal-split solution with the
static model $K=S=0$ from that same measure. At every time both models have
identical count $N_0$, mass $m$, and mean size $m/N_0$, yet their second
moments differ by
$M_2(0)(e^{3\lambda mt/2}-1)$. At a fixed positive time this difference
can be made arbitrarily large by increasing $\lambda$. Consequently these
three bulk outputs alone cannot bound the second-moment error uniformly
over the class. For each fixed $\lambda$ the second moment remains finite
at every finite time; this example has no finite-time gelation claim.
The entropy inequality provides another strictly increasing observable in
the same accepted solution class.

\subsection{A power-kernel nonstationarity theorem}

Consider a separate time-independent model,
\begin{equation}
 K(x,y)=xy^\alpha+yx^\alpha,\qquad S(x)=s x^\alpha,
 \qquad 0<\alpha\leq1,\quad s\geq0.
 \label{eq:power-kernel}
\end{equation}
Let $B_x$ be the pushforward of $b_x$ by $z\mapsto z/x$, so
$B_x((0,1))=2$ and $\int\theta B_x(\dd\theta)=1$. These fraction measures
may depend measurably on the parent. A stationary measure means a
nonnegative measure $v$ whose coagulation--fragmentation vector field is
zero for every bounded Borel test, with finite total event integrals.

\begin{theorem}[No finite-count equilibrium for power kernels]
\label{thm:power-no-equilibrium}
For \eqref{eq:power-kernel} and every such daughter law, there is no
stationary measure $v$ with finite positive count $N$ and finite positive
mass $m$. No negative moment, logarithmic moment, or moment above order
one is required.
\end{theorem}

\begin{proof}
Suppose $v$ is such a measure and put $D=M_\alpha(v)$.
The inequality $x^\alpha\leq1+x$ gives $0<D<\infty$; the total coagulation
and fragmentation event rates are $mD$ and $sD$. Testing stationarity
with one gives $(s-m)D=0$, hence $s=m$. The case $s=0$ is already excluded
because $m>0$.

At $s=m$ the integrated population vector field equals
$\int\mathcal Lf\dd v$, where
\begin{equation}
 \mathcal Lf(x)
 =x\int y^\alpha[f(x+y)-f(x)]v(\dd y)
  +mx^\alpha\int[f(\theta x)-f(x)]B_x(\dd\theta).
 \label{eq:power-generator}
\end{equation}
To verify this, symmetry combines the coagulation gains and the two
losses in the full weak field into
$\iint xy^\alpha f(x+y)\dd v\dd v-D\int xf\dd v-m\int x^\alpha f\dd v$.
Fragmentation adds $m\int x^\alpha\int f(\theta x)B_x(\dd\theta)v(\dd x)
-m\int x^\alpha f\dd v$, exactly as in \eqref{eq:power-generator}.
The total jump rate of this representation is
\begin{equation}
 q(x)=xD+2mx^\alpha,
 \qquad \int q(x)v(\dd x)=3mD.
 \label{eq:power-rate}
\end{equation}
Divide its nonnegative gain kernel by $q(x)>0$ to obtain a Markov
transition kernel $P(x,\dd z)$, and set
$\zeta(\dd x)=q(x)v(\dd x)/(3mD)$. Bounded-test stationarity is exactly
the finite-measure identity $\zeta P=\zeta$.

The singular test $g(x)=x^{-\alpha}$ need not be integrable under $v$,
but it is integrable under the rate-weighted probability:
\begin{equation}
 3mD\int g\dd\zeta=D M_{1-\alpha}(v)+2mN<\infty.
 \label{eq:power-singular-integrability}
\end{equation}
Test $\zeta P=\zeta$ first with $g\wedge k$ and let $k\uparrow\infty$.
Monotone convergence gives $\int Pg\dd\zeta=\int g\dd\zeta<\infty$.
Both gain and loss are therefore integrable, and subtraction is legitimate:
\begin{equation}
 \int\mathcal Lg\dd v=0.
 \label{eq:power-singular-balance}
\end{equation}
No jump-process construction or time-dependent population solution has
been used. Even if some daughter inverse moments are initially infinite,
the hypothetical stationary identity forces their integrated gain to be
finite by \eqref{eq:power-singular-integrability}.

Jensen for $B_x/2$ gives
$\int\theta^{-\alpha}B_x(\dd\theta)\geq2^{1+\alpha}$, so fragmentation
in \eqref{eq:power-singular-balance} contributes at least
$2mN(2^\alpha-1)$. Write its coagulation contribution as $-J$, where
\[
 J=\iint xy^\alpha[x^{-\alpha}-(x+y)^{-\alpha}]v(\dd x)v(\dd y)\geq0.
\]
We claim $J\leq\alpha mN$. Symmetrize the integrand, set $r=x+y$ and
$u=x/r$, and denote the symmetrized integrand by $rC_\alpha(u)/2$, where
\[
 C_\alpha(u)=u(1-u)^\alpha(u^{-\alpha}-1)
 +(1-u)u^\alpha((1-u)^{-\alpha}-1).
\]
Concavity of $z^\alpha$ gives
$u^{-\alpha}-1\leq\alpha(u^{-1}-1)$ and the counterpart with $1-u$.
Thus
\[
 C_\alpha(u)\leq\alpha\bigl[(1-u)^{1+\alpha}+u^{1+\alpha}\bigr]
 \leq\alpha.
\]
Integrating $r/2$ gives $mN$, proving the claim. It follows that
\begin{equation}
 \int\mathcal Lg\dd v
 \geq\bigl[2(2^\alpha-1)-\alpha\bigr]mN
 \geq(2\log2-1)\alpha mN>0,
 \label{eq:power-positive-drift}
\end{equation}
contradicting \eqref{eq:power-singular-balance}.
\end{proof}

At $\alpha=1$, $D=m$ and $\zeta=xv/m$ is the mass probability itself.
The embedded transition coagulates with probability $1/3$ by adding an
independent mass-sampled size and fragments with probability $2/3$ using
$B_x/2$. Since $\int x^{-1}\zeta(\dd x)=N/m<\infty$, stationarity would
imply
\[
 3\E(1/X)=\E[1/(X+Y)]+2\E[1/(X\Theta)]\geq4\E(1/X),
\]
a direct contradiction; here $X,Y$ are independent with law $\zeta$ and
$\Theta\mid X$ has law $B_X/2$. This illustrates why rate weighting makes
the singular test admissible.

\subsection{Failure of fractional-moment monotonicity beyond the additive kernel}

The preceding proof does not transfer the additive model's monotone
positive moments. With equal splitting, take
$v_R=\delta_1+R^{-1}\delta_R$, $R>1$, of mass two, and choose $s=m=2$.
Direct substitution gives, for $0<p<1$,
\begin{align}
 \mathcal D_p(v_R)
 &=A_p(1+R^{\alpha+p-1})
 +(1+R^{\alpha-1})[(1+R)^p-1-R^p],\label{eq:power-counterexample}\\
 A_p&=2^p+2^{2-p}-4=\frac{(2^p-2)^2}{2^p}>0.\notag
\end{align}
For verification, the two diagonal coagulation terms are
$(2^p-2)(1+R^{\alpha+p-1})$, fragmentation is
$2(2^{1-p}-1)(1+R^{\alpha+p-1})$, and the cross coagulation term is the
second term in \eqref{eq:power-counterexample}. If $p>1-\alpha$, the first
term tends to infinity, while the negative second term remains bounded:
$-1<(1+R)^p-R^p-1<0$ and $1+R^{\alpha-1}\leq2$.
Nothing is proved for $p\leq1-\alpha$: there the first term stays
bounded and this family does not settle the sign of
$\mathcal D_p(v_R)$.

In particular the half moment can have positive drift for every
$\alpha>1/2$. An explicit half-moment counterexample is $\alpha=1$, $p=1/2$, $R=64$:
\begin{equation}
 \mathcal D_{1/2}(v_{64})=27\sqrt2+2\sqrt{65}-54
 >\frac{151}{500}>0.
 \label{eq:power-half-positive}
\end{equation}
Indeed $(1414/1000)^2<2$ and $(8062/1000)^2<65$, and substitution gives
the rational lower bound. The supplement independently evaluates the
atomic double sum and certifies this inequality by rational arithmetic.
These are vector-field counterexamples. If a local mass-conserving solution
with the required initial moment derivative is separately established,
they give that derivative; no time-dependent existence assertion for
\eqref{eq:power-kernel} is made here. Since count and mass are neutral at
these data, the same sign obstruction applies to a universal initial decay
inequality for the normalized affinity.

The finite-count qualification is necessary. \citet[Proposition 4.3 and
Remark 4.4 of arXiv version 3, whose numbering we follow]{TranVan2022}
construct critical stationary
Bernstein transforms for $K(x,y)=xy$ and a constant binary breakup kernel,
which in our notation means selection $x/2$ and uniform daughter fractions.
Their mass-one convention maps to $K=2xy$, $S=mx$ at mass $m$ by
$n_t=m c_{2mt}$. Writing their stationary transform as
$F(q)=\int(1-e^{-qx})c(\dd x)$ and $G=1-F'$, their identity is
\[
 \frac{G(q)}{(1-G(q))^3}=Cq,\qquad C>0.
\]
It implies $F'(0)=1$ and
$F(q)\sim\tfrac32C^{-1/3}q^{2/3}\to\infty$: indeed $G(q)\to1$ and
$F'(q)=1-G(q)\sim(Cq)^{-1/3}$, which can be integrated using two-sided
asymptotic bounds. Monotone convergence in the Bernstein representation
then gives infinite count. These known equilibria do not contradict
\cref{thm:power-no-equilibrium}.

For $0<\alpha\leq1$, that theorem establishes only stationary exclusion.
It does not supply mass conservation, boundary convergence, logarithmic
scaling, or a final auxiliary coagulation event for time-dependent power
kernels. For $\alpha>1$, both the finite flux and the concavity estimate in
its proof can fail under count and mass assumptions alone; that regime is
outside the result.
